\chapter{第一部分：云计算平台搭建}

\section{任务1：应用容器化}

\subsection{后端镜像（多阶段构建 + 自选包）}

后端为 Flask API，提供 \lstinline!/api/ping!（健康检查）、\lstinline!/api/visits!
（写 Redis 计数，演示三层通信）、\lstinline!/api/info!、\lstinline!/api/compute!
（CPU 压测入口）四个接口。\lstinline!Dockerfile.backend! 采用两阶段构建：第一阶段
将依赖安装到独立目录，第二阶段仅复制依赖与代码，避免把 pip/编译工具带入运行镜像，
显著减小体积。在 \lstinline!requirements.txt! 中加入了自选包 \lstinline!requests!。

\begin{lstlisting}[language=Docker,caption={Dockerfile.backend（多阶段构建）},captionpos=b]
# Stage 1: 安装依赖
FROM python:3.11-slim AS builder
WORKDIR /build
COPY requirements.txt .
RUN pip install --no-cache-dir -r requirements.txt --target /build/packages
# Stage 2: 运行时镜像（更小）
FROM python:3.11-slim
WORKDIR /app
COPY --from=builder /build/packages /app/packages
COPY . .
ENV PYTHONPATH=/app/packages
EXPOSE 5000
CMD ["python", "app.py"]
\end{lstlisting}

\subsection{前端镜像（含学号姓名）}

前端为 Nginx 静态页，\lstinline!static/index.html! 中加入了本人学号与姓名用于答辩识别：
\begin{lstlisting}[language=HTML,caption={index.html 关键行},captionpos=b]
<p class="info">学号：2023112594　姓名：陶秋源　班级：计算机科学与技术04班</p>
\end{lstlisting}

\subsection{本地联调与推送 SWR}

本地执行 \lstinline!docker compose up! 启动 backend、frontend、redis 三个容器，
浏览器访问前端页面（图~\ref{fig:t1page}）点击按钮触发请求，后端日志打印
\lstinline!收到请求 GET /api/ping!、\lstinline!/api/visits! 等记录
（图~\ref{fig:t1logs}），证明「前端 Nginx $\to$ 后端 Flask $\to$ Redis」三层通信正常。

\shot{task1_frontend_page.png}{前端页面（含学号姓名）调用后端接口返回 JSON}{fig:t1page}
\shot{task1_docker_compose_logs.png}{docker compose 后端日志显示收到请求}{fig:t1logs}

镜像推送到华为云 SWR 时需注意：Docker~29.x 默认生成 OCI 格式 manifest，而 SWR 仅支持
经典 Docker 格式，故构建时使用 \lstinline!--provenance=false!。推送命令如下，推送后在
SWR 控制台可见 \lstinline!backend! 与 \lstinline!frontend! 镜像及其 Tag（图~\ref{fig:t1swr}）。

\begin{lstlisting}[language=bash,caption={推送镜像到 SWR group17},captionpos=b]
docker login -u cn-north-4@<AK> -p <SWR登录密钥> swr.cn-north-4.myhuaweicloud.com
docker tag backend:v1 swr.cn-north-4.myhuaweicloud.com/group17/backend:v1
docker push swr.cn-north-4.myhuaweicloud.com/group17/backend:v1
\end{lstlisting}

\shot{task1_swr_images.png}{SWR 控制台 group17 组织下的镜像列表（含名称与 Tag）}{fig:t1swr}

\section{任务2：CCE 集群搭建}

集群创建与节点规格已在第~1 章描述，2 个 Worker 节点均为 Ready、版本 v1.35.3
（$\geq 1.27$），验收截图见图~\ref{fig:nodes}。\texttt{KubeConfig} 从 CCE 控制台下载，
本地配置 \lstinline!kubectl! 使用其中的 external 上下文（公网 API 端点）直接操作集群。

\section{任务3：应用部署}

参考附录 A-2 模板，部署了 6 类资源：后端 Deployment（副本2、含 resources 限制、
\lstinline!envFrom! 注入 ConfigMap、\lstinline!secretKeyRef! 注入 Redis 密码）、
Redis Deployment（副本1、\lstinline!limits.memory!=256Mi$\leq$512Mi）、后端
LoadBalancer Service（注解 \lstinline!kubernetes.io/elb.class: union! +
\lstinline!elb.autocreate! 自动创建公网 ELB）、Redis ClusterIP Service、ConfigMap
（\lstinline!REDIS_HOST=redis-svc!）、Secret（密码 base64 编码，无明文）。

\begin{lstlisting}[language=yaml,caption={后端 Deployment 关键片段},captionpos=b]
spec:
  replicas: 2
  template:
    spec:
      containers:
      - name: backend
        image: swr.cn-north-4.myhuaweicloud.com/group17/backend:v1
        resources:
          requests: { cpu: "100m", memory: "128Mi" }
          limits:   { cpu: "500m", memory: "512Mi" }
        envFrom:
        - configMapRef: { name: backend-config }   # 注入 Redis 地址
        env:
        - name: REDIS_PASSWORD
          valueFrom:
            secretKeyRef: { name: redis-secret, key: password }  # 注入密码
\end{lstlisting}

部署后所有 Pod 均为 Running（图~\ref{fig:t3pods}），后端 Service 自动绑定华为云公网
ELB，分配公网 IP \textbf{1.94.238.112}（图~\ref{fig:t3svc}）。浏览器访问
\lstinline!http://1.94.238.112/api/ping! 返回 \lstinline!{"status":"ok"}!
（图~\ref{fig:t3ping}）；连续访问 \lstinline!/api/visits! 计数依次递增且
\lstinline!served_by! 在两个后端 Pod 间轮换，验证了 ELB 负载均衡、ConfigMap/Secret
注入与三层通信全部正常。ConfigMap 与 Secret 内容见图~\ref{fig:t3cmsecret}，Secret
的 password 字段为 base64（\lstinline!Z3JvdXAxNy4=!），无明文。

\shot{task3_get_pods.png}{kubectl get pods：backend×2、redis、frontend 全部 Running}{fig:t3pods}
\shot{task3_get_svc_elb.png}{后端 Service 为 LoadBalancer，公网 IP 1.94.238.112}{fig:t3svc}
\shot{task3_api_ping.png}{浏览器访问 ELB 公网 IP，/api/ping 返回 status:ok}{fig:t3ping}
\shot{task3_configmap_secret.png}{ConfigMap 注入 Redis 地址；Secret 密码 base64 无明文}{fig:t3cmsecret}

\section{任务4：持久化存储}

为 Redis 配置 PVC（\lstinline!storageClassName: csi-disk!，即华为云 EVS 云硬盘，
申请 10\,Gi），并修改 Redis Deployment 将 \lstinline!/data! 挂载到该 PVC；同时 Redis
以 \lstinline!--appendonly yes! 开启 AOF 持久化。由于 EVS 为 \lstinline!ReadWriteOnce!，
Deployment 更新策略设为 \lstinline!Recreate! 以避免新旧 Pod 同时挂载同一磁盘。

PVC 创建后状态为 \textbf{Bound}（图~\ref{fig:t4pvc}）。持久化验证流程：向 Redis 写入
\lstinline!SET testkey "hello"!（图~\ref{fig:t4write}）$\to$ \lstinline!kubectl delete pod!
删除 Redis Pod 触发重建（图~\ref{fig:t4delete}）$\to$ 新 Pod Running 后再次
\lstinline!GET testkey! 仍返回 \lstinline!"hello"!（图~\ref{fig:t4persist}），证明数据
随云硬盘持久化，未随 Pod 销毁而丢失。

\shot{task4_pvc_bound.png}{kubectl get pvc：redis-data-pvc 状态 Bound（10Gi csi-disk）}{fig:t4pvc}
\shot[0.8]{task4_redis_write.png}{写入测试数据 SET testkey hello}{fig:t4write}
\shot[0.8]{task4_pod_delete.png}{删除 Redis Pod，触发自动重建}{fig:t4delete}
\shot[0.8]{task4_redis_persist.png}{重建后 GET testkey 仍返回 hello（数据持久化）}{fig:t4persist}

\section{任务5：ConfigMap Volume 挂载}

将前端 Nginx 的反向代理配置（\lstinline!default.conf!）改为以 ConfigMap 的 Volume
形式挂载到 \lstinline!/etc/nginx/conf.d/default.conf!（\lstinline!subPath! 挂载），
其中 upstream 指向后端 Service \lstinline!backend-internal:5000!。

随后将 ConfigMap 中后端端口由 5000 改为 5001 并 \lstinline!kubectl apply!。注意
Kubernetes 已知限制：使用 \lstinline!subPath! 挂载时 ConfigMap 更新\textbf{不会}自动
同步到 Pod 内，故需删除前端 Pod 让 Deployment 重建。重建后 \lstinline!exec! 进入前端
Pod 执行 \lstinline!cat default.conf!，可见端口已变为 5001（图~\ref{fig:t5after}），
证明挂载文件随 ConfigMap 更新而生效。

\shot[0.8]{task5_nginx_before_5000.png}{修改前：挂载配置 proxy\_pass 指向 :5000}{fig:t5before}
\shot[0.8]{task5_nginx_after_5001.png}{修改并重建后：挂载配置已更新为 :5001}{fig:t5after}

\subsection{Volume 挂载与 envFrom 两种方式的适用场景}

\lstinline!envFrom!（环境变量注入）适合\textbf{少量、扁平的 key-value 配置}（如开关、
地址、端口），应用通过 \lstinline!os.environ! 读取，简单直接，但环境变量在容器启动时
固定，更新需重建 Pod，且不适合大段结构化文本。\textbf{Volume 挂载}则适合\textbf{完整
配置文件}（如 \lstinline!nginx.conf!、\lstinline!application.yaml!）：以文件形式呈现，
保留原有格式，且在非 \lstinline!subPath! 挂载时支持 kubelet 周期性热更新（约分钟级）。
本设计后端的 Redis 地址用 \lstinline!envFrom! 注入，前端的整段 Nginx 配置用 Volume
挂载，正是按"短变量用环境变量、长文件用卷"的原则选择。

\section{任务6：HPA 弹性伸缩}

确认 \lstinline!kubectl top nodes! 有数据后（metrics-server 就绪），为后端创建 HPA：
\lstinline!minReplicas=1!、\lstinline!maxReplicas=4!、CPU 目标利用率 60\%。随后用压测
脚本对 ELB 公网 IP 的 CPU 密集端点 \lstinline!/api/compute! 发起约 120 并发的持续请求。

\shot[0.8]{task6_top_nodes.png}{kubectl top nodes：metrics-server 提供资源指标}{fig:t6top}
\shot[0.8]{task6_hpa_apply.png}{创建 HPA：min1/max4/CPU 60\%}{fig:t6hpa}

压测开始后约 1--2 分钟，\lstinline!kubectl get pods -w! 监控窗口可见后端 Pod 数量从 1
增加到 2 乃至更多（图~\ref{fig:t6up}）；停止压测约 5 分钟后，Pod 数缩回 1
（图~\ref{fig:t6down}）。

\shot{task6_scale_up.png}{压测中后端 Pod 从 1 扩容到多个}{fig:t6up}
\shot{task6_scale_down.png}{停止压测后 Pod 数缩回 1}{fig:t6down}

\subsection{结果分析：扩容延迟、冷却时间与降本价值}

\textbf{① 扩容延迟的原因。} 从 CPU 飙升到新 Pod 就绪存在分钟级延迟，源于多段串联：
metrics-server 默认每 15\,s 采集一次指标、HPA 控制器默认每 15\,s 评估一次、计算出目标副本
后还需经历调度、镜像拉取、容器启动与就绪探针通过。这些周期叠加导致扩容并非"瞬时"。

\textbf{② 冷却时间（稳定窗口）的意义。} HPA 缩容默认有 300\,s 的稳定窗口
（\lstinline!stabilizationWindowSeconds!），即负载下降后要持续一段时间才真正缩容。
其作用是\textbf{防止抖动}：避免负载在阈值附近波动时副本数频繁增减（thrashing），
保证服务稳定、减少 Pod 反复创建销毁的开销。

\textbf{③ HPA 对降本的价值。} HPA 让副本数随真实负载自动伸缩——低峰期仅保留
\lstinline!minReplicas=1! 节省 CPU/内存与节点开销，高峰期自动扩容保障 SLA。相比按
峰值固定容量部署，弹性伸缩显著提升资源利用率、降低云上成本，是云计算"按需付费、
弹性供给"理念的直接体现。
